% Generated draft: requires mathematical and source review.
\documentclass[11pt]{article}
\usepackage[margin=25mm]{geometry}
\usepackage{fontspec}
\setmainfont{Latin Modern Roman}
\usepackage{amsmath,amssymb,mathtools,unicode-math}
\setmathfont{Latin Modern Math}
\usepackage{enumitem,needspace,xurl,hyperref}
\setlength{\parindent}{0pt}
\setlength{\parskip}{5pt}
\setlength{\emergencystretch}{3em}
\newcommand{\N}{\mathbb{N}}
\newcommand{\Z}{\mathbb{Z}}
\newcommand{\Q}{\mathbb{Q}}
\newcommand{\R}{\mathbb{R}}
\newcommand{\C}{\mathbb{C}}
\newcommand{\F}{\mathbb{F}}
\newcommand{\A}{\mathbb{A}}
\newcommand{\PP}{\mathbb{P}}
\newcommand{\E}{\mathbb{E}}
\newcommand{\eps}{\varepsilon}
\newcommand{\dd}{\,\mathrm{d}}
\begin{document}
\section*{UnsolvedMath ID: 10000075}
\textbf{Source problem code:} catalogue entry
\subsection{Definitions and mathematical statement}
Fix \(0<C<1/4\). The source considers uniform rooted quadrangulations with \(N\) faces and genus \(CN\). Since genus is integral, a precise asymptotic model must choose admissible integers \(g_N\) with \(g_N/N\to C\), or restrict to a subsequence on which \(CN\) is integral. The map class and root convention are not specified.

For rooted maps, local convergence means convergence in distribution of every fixed-radius rooted neighborhood. The proposed limit is the source's stochastic hyperbolic infinite quadrangulation (SHIQ), said to arise from a labeled supercritical Galton--Watson tree. Prove the intended rooted local convergence, the asserted mean degree
\[
 \frac{6}{1-4C},
\]
and divergence of local injectivity radius. A conventional probabilistic version of the last claim is that, for every fixed \(R\), the probability that the radius-\(R\) neighborhood detects nontrivial surface topology tends to zero. Injectivity here concerns the universal cover of the surface, not the graph-theoretic tree cover, which would incorrectly discard face cycles.

The excerpt omits the offspring law, label process, map construction and limiting root distribution, so the named SHIQ object is not uniquely defined by the phrase ``supercritical Galton--Watson tree''. The meaning of mean degree is also material. For a closed orientable quadrangulation, Euler's formula gives \(E=2N\), \(V=N+2-2g_N\), so the empirical mean at a uniformly sampled vertex tends to \(4/(1-2C)\). The displayed \(6/(1-4C)\) is instead the corresponding Euler mean for triangulations with \(N\) faces. An edge-rooted degree law need not equal the uniform-vertex mean. Hence the map type, root and asserted mean may be mixed and require original-source verification; neither formula is silently substituted for the source prediction.

The registry labels the item solved, but the summary says partial/major or largely solved, asks to verify the citation, and supplies the placeholder ``arXiv:1904.xxxx''. It attributes high-genus local-limit work to Budzinski--Curien and later authors without a precise theorem matching this exact quadrangulation, tree and mean-degree statement. The reported progress is retained as an attribution rather than a verified resolution.
\subsection{Short English statement}
Identify the local limit of uniform quadrangulations whose genus grows linearly with face count, and verify the proposed degree and injectivity-radius claims. The limiting tree construction is missing, and the stated mean-degree formula may mix triangulation and quadrangulation or rooting conventions.
\subsection{Sources}
\begin{itemize}
\item Local registry: \url{data/unsolvedmath/problems.json}, permanent UnsolvedMath ID \(10000075\), source code catalogue entry; supplied statement, background, and research summary, updated 2026-08-06. Attributed to Itai Benjamini, 2012.
\item Benjamini, \emph{Euclidean vs. Graph Metric}, section 8.1, unnumbered SHIQ conjecture, PDF page 11: \url{https://www.wisdom.weizmann.ac.il/~itai/erd100.pdf#page=11}. Preserved copy: \url{https://arquivo.pt/noFrame/replay/20201231041538id_/http://www.wisdom.weizmann.ac.il/~itai/erd100.pdf}. Supplied extraction checked against rendered PDF, accessed 2026-07-29.
\item Budzinski--Curien, ``Random maps with large genus'', and ``Hyperbolic random maps'', as described by the summary. Its arXiv:1904.xxxx is an incomplete placeholder, not a usable identifier; exact citations and hypotheses are absent.
\end{itemize}
\end{document}
